% capitulo3.tex
\chapter{Hipótesis}

% ============================================================
% SECCIÓN 3.1 — FORMULACIÓN DE LAS HIPÓTESIS
% ============================================================
\section{Formulación de las Hipótesis de Investigación}

Las hipótesis que guían el presente trabajo se formulan en función de los cambios normativos
específicos identificados entre la NTC-DCEA 2017 y la NTC-DCEA 2023. Cada hipótesis es verificable
mediante los algoritmos paramétricos implementados y los cuatro barridos de cálculo ejecutados con
el perfil IR~305$\times$74.5 como caso semilla.

% ============================================================
% SECCIÓN 3.2 — HIPÓTESIS 1
% ============================================================
\section{Hipótesis 1: Espesor de la Placa Base}

Los cambios introducidos por la NTC-DCEA 2023 en el diseño de placas base bajo flexocompresión
con gran excentricidad no producen variaciones significativas en el espesor requerido $t_{req}$
de la placa, dado que ambas normas emplean el mismo modelo mecánico del bloque equivalente de
compresión de Whitney para determinar la profundidad del bloque $Y$, la fuerza de tensión
$T_{ua}$ en los pernos de anclaje y, en consecuencia, el momento flector de diseño de la placa.
Las diferencias en $t_{req}$, si existen, se esperan únicamente bajo muy alta excentricidad,
donde la NTC-2023 exige verificar el mayor momento entre el lado de tensión y el lado de
compresión, lo que puede añadir hasta 2~mm al espesor calculado respecto a la norma anterior.

% ============================================================
% SECCIÓN 3.3 — HIPÓTESIS 2
% ============================================================
\section{Hipótesis 2: Resistencia al Cono de Concreto}

La incorporación de los factores de confinamiento $R_{et}$ y $R_{ev}$ por refuerzo transversal
del dado en la NTC-DCEA 2023 produce un incremento sustancial en la resistencia nominal al
arrancamiento del cono de concreto respecto a la metodología de la NTC-DCEA 2017. Este
incremento es función creciente de la profundidad de empotramiento $h_{ef}$ y de la densidad
del refuerzo transversal ---cuantificada por el número de estribos interceptados por el plano
del cono hipotético, $N_{est} = \lfloor h_{ef}/s_{est} \rfloor$--- y representa una reducción
del conservadurismo excesivo del diseño anterior sin comprometer los márgenes de seguridad
reglamentarios.

% ============================================================
% SECCIÓN 3.4 — HIPÓTESIS 3
% ============================================================
\section{Hipótesis 3: Distancia Mínima al Borde}

El incremento del 16.7\,\% en la distancia mínima al borde libre requerida por la NTC-DCEA 2023
para anclas de alta resistencia ($6\,d_o \to 7\,d_o$) tiene implicaciones directas y
cuantificables en las dimensiones mínimas en planta del dado de concreto, constituyendo la
modificación con mayor impacto geométrico sobre la subestructura para diámetros de ancla
superiores a $3/4"$.
